% lemma_run.tex — the local law along a run of quadratic residues and the descent formula for signature densities
% (first solver, 2026-08-19).  For the section on the neighbour-graph components (k=-1).  Uses \Fp, amsthm; labels lem:run, prop:descent.

\begin{lemma}[equidistribution along a run]\label{lem:run}
Fix $k\ge2$ and let $C_k$ be the affine curve
$\{(u_1,\dots,u_k):\ u_{i+1}^2=u_i^2+1\ (1\le i<k)\}$, i.e.\ the multiquadratic cover
$\Fp(t)\bigl(\sqrt t,\sqrt{t+1},\dots,\sqrt{t+k-1}\bigr)$ of the $t$-line ($t=u_1^2$).  It is absolutely irreducible of genus
$1+2^{k-2}(k-3)$, and for every $\mathbf h=(h_1,\dots,h_k)\in\Fp^k\setminus\{0\}$ the function $\sum_ih_iu_i$ is non-constant on $C_k$.
Consequently, for all intervals $I_1,\dots,I_k\subseteq(-p/2,p/2)$ and all $\varepsilon_0,\varepsilon_k\in\{\pm1\}$,
\[
 \#\Bigl\{(u_i)\in C_k(\Fp):\ X(u_i)\in I_i\ \forall i,\ \chi(u_1^2-1)=\varepsilon_0,\ \chi(u_k^2+1)=\varepsilon_k\Bigr\}
 =\frac p4\prod_{i=1}^k\frac{|I_i|}p+O_k\bigl(\sqrt p\log^kp\bigr),
\]
where $\chi$ is the quadratic character and $X$ the centred representative.  In particular the positions $X(u_i)/p$ are asymptotically
independent and uniform on $(-\frac12,\frac12)$, and independent of the two Legendre conditions that make the run maximal.
\end{lemma}
\begin{proof}
Irreducibility: the $k$ quadratic extensions are independent because the polynomials $t,t+1,\dots,t+k-1$ are pairwise coprime and squarefree
(for $p>k$), so the Galois group over $\overline\Fp(t)$ is $(\Z/2)^k$ and the cover is connected.  Genus: the cover of degree $2^k$ is
ramified exactly over the $k$ finite points $t=0,-1,\dots,-(k-1)$ and over $\infty$, with inertia $\Z/2$ at each (at $\infty$ the inertia is
generated by the simultaneous flip of all square roots, each $\sqrt{t+i}$ having odd polar order); each such point has $2^{k-1}$ preimages
with $e=2$, so $2g-2=2^k(-2)+(k+1)2^{k-1}$, i.e.\ $g=1+2^{k-2}(k-3)$.  Non-constancy: the group $(\Z/2)^k$ acts on $C_k$ by
$u_i\mapsto\pm u_i$ independently; if $\sum_ih_iu_i$ were constant, averaging over the subgroup flipping only $u_j$ would give $h_ju_j$
constant, hence $h_j=0$ for every $j$.  Now expand the interval indicators in Fourier series (or Selberg polynomials) of $\ell^1$-norm
$O(\log p)$ each and the two Legendre conditions as $\frac12(1\pm\chi)$; the main term is the product of the volumes times
$|C_k(\Fp)|/4=p/4+O_k(\sqrt p)$, and every other term is a mixed character sum
$\sum_{C_k(\Fp)}\chi\bigl((u_1^2-1)^{a}(u_k^2+1)^{b}\bigr)e\bigl(\sum_ih_iu_i/p\bigr)$ with $(a,b,\mathbf h)\ne0$, which is $O_k(\sqrt p)$ by the
Bombieri--Perel'muter bound: the additive part is non-constant when $\mathbf h\ne0$, and for $\mathbf h=0$ the Kummer sheaf is geometrically
non-trivial because $u_1^2-1=t-1$ and $u_k^2+1=t+k$ have odd order at their zeros on the $t$-line and their zeros are unramified in $C_k$.
\end{proof}

\begin{proposition}[signature densities are descent statistics]\label{prop:descent}
Let a component of the neighbour graph be a maximal run $t,t+1,\dots,t+k-1$ of quadratic residues, and for $1\le i\le k-1$ let
$\xi_i\in\{0,1\}$ record whether the $H(1)$-pair carried by the edge $(t+i-1,t+i)$ shares its centre.  Then, as $p\to\infty$,
\[
 \Pr\bigl[(\xi_1,\dots,\xi_{k-1})=\varepsilon\bigr]=\frac{D_k(S_\varepsilon)}{k!}+O_k\bigl(p^{-1/2}\log^kp\bigr),\qquad
 S_\varepsilon=\{i:\varepsilon_i=1\},
\]
where $D_k(S)$ is the number of permutations of $\{1,\dots,k\}$ with descent set exactly $S$.  The types of the vertices are determined by
$\varepsilon$: the vertex between the edges $i-1$ and $i$ is a good group $(4,8,4)$ iff $(\xi_{i-1},\xi_i)=(0,1)$, a $(2,6,6,2)$-group iff
$(1,0)$, and a seven-group otherwise; the end vertices are $(4,2,2)$ iff $\xi_1=1$ (left), resp.\ $\xi_{k-1}=0$ (right), and $(3,3,1,1)$ otherwise.
\end{proposition}
\begin{proof}
Write $u_i$ for a square root of $t+i-1$.  The $H(1)$-pair of the edge $(t+i-1,t+i)$ is $\{a,1/a\}$ with
$a=u_i+u_{i+1}$, $1/a=u_{i+1}-u_i$ (indeed $a\cdot(1/a)=u_{i+1}^2-u_i^2=1$ and $a-1/a=2u_i$ is the residue $d$ of the group), and the pair
shares its centre iff $X(a)X(1/a)<0$; this is independent of the choice of signs of the $u_i$'s.  Put $v_i=X(u_i)/p$.  Since the sign of the
centred representative of $x$ equals the sign of $\sin2\pi x$ and
$\sin2\pi(v+w)\,\sin2\pi(w-v)=\tfrac12\bigl(\cos4\pi v-\cos4\pi w\bigr)$, we get the exact identity
\[
 \xi_i=1\iff\cos4\pi v_i<\cos4\pi v_{i+1}\iff g(v_i)>g(v_{i+1}),\qquad g(x):=\min\bigl(|x|,\tfrac12-|x|\bigr).
\]
By Lemma~\ref{lem:run} the $v_i$ are asymptotically independent and uniform, hence the $g(v_i)$ are independent and uniform on $(0,\frac14)$
(the map $g$ is four-to-one and piecewise linear), so $(\xi_i)$ is the descent pattern of $k$ i.i.d.\ uniform variables, whose law is
$D_k(S)/k!$.  The type dictionary is Lemma~\ref{lem:partners} together with the fact that the pair on an edge is seen as an
$H(1)$-pair by its left vertex and as (the $R$-image of) a $\tau$-pair by its right vertex, which exchanges $A\cup D$ with $B\cup C$.
\end{proof}
